\section{Introduction}\label{main:introduction}

Metric $k$-median asks for a small set of facilities serving a finite set
of clients.  The input consists of finite sets $D$ of clients and $F$ of
available facilities, distances given by a finite rational metric on their
locations, and a positive integer $k$.  For a nonempty $S\subseteq F$ set
\[
 \operatorname{cost}(S)=\sum_{p\in D}d(p,S),\qquad
 d(p,S)=\min_{s\in S}d(p,s),
 \qquad
 \operatorname{opt}_k=\min_{1\le|S|\le k}\operatorname{cost}(S).
\]
A feasible solution opens at most $k$ facilities.  Clients and candidate
facilities are specified separately; they may share locations.
We write $n=|D|$ and $N=|D|+|F|$.

This paper develops two tools for respecting that budget.  The first
recovers a good solution from an anchor that represents all but a small
number of comparison clusters accurately.  Its running time is single-exponential
in the number of exceptions, making logarithmically many exceptions
tractable.  The second gives an elementary proof of bounded-price
strictness for one fixed execution of a facility-opening process.
The recovery and execution arguments combine into a randomized
approximation below $2$ on arbitrary finite rational metrics.

\subsection{Recovery from an anchor}

Fix a comparison solution $O\subseteq F$ of size $h$ and cost $P>0$.
Assign every client to a nearest member of $O$, using fixed priorities
for ties, and write $C_i$ for the clients assigned to $i\in O$ and
$P_i=\sum_{p\in C_i}d(p,i)$.  Clusters may be empty.
An \emph{anchor} is a supplied feasible set $S\subseteq F$ of the same
size $h$.  Partition $O$ into good and bad centers.  For each good
center $i$, a \emph{proxy} is a facility $\phi_i\in S$, with the proxies
pairwise distinct.  Put $P_g=\sum_{i\text{ good}}P_i$.
The recovery promise bounds the total proxy-assignment cost by
$P_g+\mu P$, where $\mu$ is a fixed accuracy parameter.  The comparison
solution, its partition and the proxy map certify this promise; the
algorithm is not given any of them.

The precise recovery theorem uses a normalized representation: clients
and facilities have disjoint indices, distances between distinct indices
are positive integers, and their maximum is bounded by a fixed polynomial
in $N$.  The fixed polynomial and accuracy parameters may affect the
running-time degree.

\begin{theorem}[Refined recovery]\label{ref:recovery}
Let $D$ and $F$ be disjoint indexed sets in a metric whose distances
between distinct indices are positive integers bounded by a fixed
polynomial in $N=|D|+|F|$.
Fix $\varepsilon>0$ and $0<L_0<\infty$, and choose a fixed rational parameter
\[
 0<\mu\le \min\{1/4,\varepsilon/12\}.
\]
Let $S\subseteq F$ have cardinality $h\ge1$.  Suppose there exists a comparison
solution $O\subseteq F$ of cardinality $h$ and cost $P>0$, with a partition into
good and bad centers and distinct proxies $\phi_i\in S$ for the good centers,
such that
\begin{equation}\label{ref:proxy-hypothesis}
 \sum_{i\ \mathrm{good}}\sum_{p\in C_i}d(p,\phi_i)
 \le P_g+\mu P,
 \qquad m:=|\{i:i\ \mathrm{bad}\}|\le L_0\log N.
\end{equation}
Given $S$, $\mu$, $L_0$, and a failure parameter $\zeta\in(0,1)$, a randomized
algorithm returns $\widehat S\subseteq F$ with
\[
 |\widehat S|\le h
 \quad\text{always},\qquad
 \operatorname{cost}(\widehat S)
 \le (1+2/e+\varepsilon)P
 \quad\text{with probability at least }1-\zeta.
\]
Its running time is polynomial in $N$ and $\log(1/\zeta)$ for the fixed
parameters.  The algorithm is not given $O$, $P$, the partition, or its proxies.
\end{theorem}

The running time is single-exponential in the allowed number of bad
centers.  In particular, Theorem~\ref{comp:recovery} gives factor $1.97$
with proxy error $P/100$ and at most $M_0$ bad centers in time
\[
 \exp(O(M_0))\operatorname{poly}(N)\log(2/\zeta).
\]
Local optimality of the anchor is not part of either recovery promise.
Section~\ref{loc:section} proves one way to produce such proxies from
local optimality and separation.  Declaring all centers bad also gives
the promise when $h=O(\log N)$.

The normalized promise matters.  Theorem~\ref{ref:recovery} does not
assert that the proxies of an arbitrary original metric retain their
accuracy after distances are capped and rounded.  The general-metric
application below instead constructs its anchors \emph{after}
normalization and proves a separate transfer of the final cost bound.

\subsection{A compatible execution and its global application}

The other tool concerns client budgets produced by the logarithmic-surplus
construction of Cohen-Addad, Grandoni, Lee, Schwiegelshohn and Svensson
\cite{CGLSS2026}.  At target cardinality $h$, that construction uses
$h$ paid regular facilities and $O(\log N)$ free copies, apart from a
small-price endpoint handled separately.  A free copy is a labeled leaf
attached to an available facility.  Replacing each copy by its underlying
facility produces a real solution without increasing connection cost.

The analysis must use one final price, one set of final copy lengths,
and the budgets of one actual replay.  Let $\lambda$ be its regular
opening payment and $\alpha_p$ the final budget of client $p$.
For a client set $C$ served by a comparison facility $i$, put
$P_C=\sum_{p\in C}d(i,p)$.  The usual dual comparison charges at most
$\lambda+2P_C$ to this set.  Lemma~\ref{str:strictness} proves that,
for each fixed $T>0$, the condition $\lambda\le TP_C$ improves this to
\[
 \sum_{p\in C}\alpha_p
 \le \lambda+(2-c(T)/2)P_C,
 \qquad c(T)>0,
\]
with constants and execution hypotheses stated there.  Crucially, the
entire opening payment $\lambda$ remains available.  Summing such
inequalities can therefore improve connection cost while still paying
for all $h$ regular openings.

Bounded-price strictness itself is established prior work.  The
factor-revealing analysis of Cohen-Addad, Grandoni, Lee and
Schwiegelshohn \cite[Appendix~B]{CGLS2026} gives a stronger numerical
gap, and its finite dual also applies when the phase variables are not
ordered and the first-connection constraint is omitted.  Our
suffix-maximum and complementary-slack argument is an alternative
self-contained proof.  Its role here is to connect that phenomenon to
the final valid-sequence replay, including free openings, witness
supersets and simultaneous phase updates.  Section~\ref{imp:section}
proves the local replay guarantees and supplies the free-opening case
needed in the final-dual argument.

The two tools give the following application on the original input
metric.  Its factor is an absolute constant; its significance here is
the way the recovery and payment arguments enforce the budget together.
The separate companion \emph{The approximation threshold for metric
$k$-median} \cite[Theorem~1.1]{OpenAIThreshold2026} proves a deterministic
$(1+2/e+\varepsilon)$ approximation for every fixed $\varepsilon>0$
on this general facility model.  The recovery, execution and global
application proofs below are independent of that graph-rounding argument.

\begin{theorem}[Global application]\label{main:theorem}
There is an absolute constant $\sigma>0$ with the following property.
For every fixed $a>0$, a randomized polynomial-time algorithm for metric
$k$-median always returns a set $S\subseteq F$ with $|S|\le k$ and satisfies
\[
 \Pr\!\left\{\operatorname{cost}(S)
       \le(2-\sigma)\operatorname{opt}_k\right\}
 \ge 1-(|D|+|F|+2)^{-a}.
\]
The same algorithm can be chosen to satisfy
$\mathbb E\operatorname{cost}(S)\le(2-\sigma)\operatorname{opt}_k$.
The running time is polynomial in the bit length of the finite rational
metric input. Its degree may depend on the fixed constants.
\end{theorem}

\subsection{Prior work and the relation to these methods}

The development of metric $k$-median approximations includes LP rounding
\cite{CharikarGuhaTardosShmoys2002}, primal-dual and Lagrangian methods
\cite{JainVazirani2001}, dual-fitting greedy facility location
\cite{JMS2002,JainEtAl2003}, and bounded-exchange local search
\cite{Arya2004}.  These methods expose different aspects of the budget:
a facility-location price controls a weighted opening cost, whereas
$k$-median requires a cardinality bound on every output.  Li and
Svensson's pseudo-approximation reduction \cite{LiSvensson2016} and
subsequent bi-point rounding \cite{ByrkaEtAl2017,GowdaEtAl2023}
provided major advances in translating relaxed solutions into feasible
ones.  More recent iterative randomized rounding gives another route
through this issue \cite{ByrkaEtAl2026}.

The closest algorithmic predecessor here is the $(2+\varepsilon)$
approximation of Cohen-Addad et al.\ \cite{CGLSS2026}, announced at
STOC 2025; we use the full version of 19 May 2026.  Its construction
combines logarithmic facility surplus with recovery for instances
whose optimum is sensitive to deleting a center.  We retain that
architecture and use its valid-sequence merging construction as an
external input.  Sections~\ref{imp:section}--\ref{sur:section} state
its exact interface, justify our scale and precision choices, and
prove the replay properties used in the application.

The recovery track is also a substantial adaptation of that work.
Its Section~4.1.1 obtains almost exact proxies outside exceptional
clusters, drawing on the structural local-search analysis of
Cohen-Addad and Schwiegelshohn \cite{CohenAddadSchwiegelshohn2017}.
Its Sections~4.2--4.4 sample leaders by distance, guess facility balls,
attach leaf dummies and remove original centers; Section~4.6 completes
these choices by partition-constrained submodular optimization.
We prove the proxy statement under an explicit scale hypothesis,
interleave the leader and removal steps, and amortize their radius
guesses through disjoint sets of clients.  An exact subset dynamic
program handles removal destinations without a sampled ball, while a
submodular surrogate handles destinations represented by proxies or
balls.  These changes yield the single-exponential exceptional-count
bound used here.

The factor $1+2/e$ in recovery has an established FPT antecedent in
Cohen-Addad, Gupta, Kumar, Lee and Li \cite{CGKLL2019}.  Their
leader/ball construction and submodular savings formulation explain
the appearance of the $1-1/e$ gain.  The continuous-greedy method and
partition rounding are due to Vondr\'ak and to C\u{a}linescu, Chekuri,
P\'al and Vondr\'ak \cite{Vondrak2008,CCPV2011}; we give the finite
categorical implementation needed for our surrogate.  Dai, Li and
Peng \cite{DaiLiPeng2026} give a recent related FPT construction.
The $k^{O(k)}$ dependence in these FPT bounds does not become polynomial
merely by setting $k=O(\log N)$.  The distinction addressed by our
recovery theorem is the dependence on the number of bad comparison
clusters, under its stated proxy promise.

\subsection{Proof overview}

The recovery proof begins with a real anchor $S$ and uses a comparison
set $O$ only to identify a favorable random history.  When local search
is used to create the anchor, weighted swaps protect the distinct good
proxies and bound the cost carried by the bad clusters.  A geometric
threshold argument then bounds their number; see
Section~\ref{loc:section}.

For recovery itself, a state consists of surviving members of $S$ and
leaf dummies attached to sampled clients.  Each step samples a client
in proportion to its distance from this reference set and then either
adds a guessed ball around it or removes its currently serving original
facility.  A population radius, computable from the input and guessed
comparison cost, controls the radius menu.  Large menu indices can
occur only at disjoint client sets whose total comparison cost pays
for them.  Their total probability loss is thus $\exp(-O(m))$, even
when the original distance range is large.  Section~\ref{sam:section}
proves this amortization and sums the likelihoods of all favorable
adaptive histories.  Saving every prefix lets the algorithm use a
useful stopping state without knowing the comparison partition.

Section~\ref{comp:section} turns each prefix into a real solution.
Suppose $a$ bad clusters have balls and a set $Q$ of original centers has
already been removed.  The analysis identifies $m-|Q|$ further removals and
at most $m-a$ destinations at comparison centers without balls.
A subset dynamic program optimizes those destinations.  For the others,
choosing one facility per ball maximizes a monotone submodular saving,
including the cost of the remaining removals.  The construction adds
at most $m$ real facilities after removing $m$ original ones, so every
candidate respects the original budget.  The comparison centers are
used to bound the optimized surrogate, never as algorithm advice.

The sharper analysis in Section~\ref{ref:section} uses two successive
stopping rules for the same sampling experiment.  At the first stop,
reference cost is small both on nearby uncovered clients and on clients
served by original centers that still need removal.  The guide then
adds balls without making further removals.  The latter client set
can only shrink.  Its fractions measured at the first stop distinguish
the clusters needing another ball from those already cheap enough.
Keeping these fractions in the completion estimate gives the
$1+2/e+\varepsilon$ guarantee.

The execution track begins in Section~\ref{imp:section}.  We fix the
surplus construction's final price and copy lengths, then replay only
its final valid sequences.  One client-budget potential pays regular
openings, and exact maximality of each phase gives a no-overbidding
inequality at actual atomic boundaries.  The final-dual proof treats
regular and free openings separately.  All three conclusions therefore
refer to the same budget vector.  Section~\ref{str:section} orders
client removals and forms a suffix maximum.  Two complementary slack
estimates show that near equality in every budget inequality would
force incompatible early and late behavior.  This gives the bounded-price
gap uniformly in the number of clients.

Finally, Sections~\ref{sur:section}--\ref{asm:section} join the tracks.
Let $L=O(\log N)$ bound the surplus.  If lowering the budget from $k$
to $k-2L$ changes optimum cost little, the target $k-L$ leaves room
for the surplus; anchoring and bounded-price strictness improve the
cost below $2$.  Otherwise a nearby cardinality has expensive
single-center deletions.  Its optimum is sufficiently separated for
anchoring to leave only $O(L)$ bad clusters, and recovery is polynomial.
For $k\le2L$ all clusters may be declared bad.  The algorithm computes
the candidates from both cases and keeps the cheapest, without testing
optimum costs.  The normalization of Section~\ref{pre:section} and a
bounded-cost local-search fallback give both the probability and the
expectation conclusions on the original metric.
